\DocumentMetadata{
  pdfversion=2.0,
  pdfstandard=ua-2,
  lang=en-US,
  tagging=on,
  tagging-setup={math/setup=mathml-SE,tabsorder=structure}
}
\documentclass[11pt,letterpaper]{article}
\usepackage[margin=0.68in,includefoot]{geometry}
\usepackage{newcomputermodern}
\usepackage{amsmath,amscd,mathtools}
\usepackage{tikz,adjustbox}
\providecommand{\square}{\mdlgwhtsquare}
\usepackage[hidelinks]{hyperref}
\hypersetup{
  pdftitle={Analysis PhD Exam, September 2000},
  pdfauthor={University of Florida Department of Mathematics},
  pdfsubject={Graduate mathematics examination},
  pdfdisplaydoctitle=true,
  bookmarksopen=true
}
\setcounter{secnumdepth}{0}
\setlength{\parindent}{0pt}
\setlength{\parskip}{0.28em}
\setlength{\emergencystretch}{3em}
\setlength{\leftmargini}{2.15em}
\setlength{\leftmarginii}{2.65em}
\setlength{\leftmarginiii}{3em}
\renewcommand{\labelenumi}{\textbf{\theenumi.}}
\pagestyle{empty}
\raggedbottom
\allowdisplaybreaks
\newenvironment{examproblems}[1]{%
  \begin{enumerate}%
  \setcounter{enumi}{#1}%
  \setlength{\topsep}{0.35em}%
  \setlength{\partopsep}{0pt}%
  \setlength{\itemsep}{0.48em}%
  \setlength{\parsep}{0.18em}%
}{\end{enumerate}}
\newenvironment{examsubparts}{%
  \begin{enumerate}%
  \setlength{\topsep}{0.18em}%
  \setlength{\partopsep}{0pt}%
  \setlength{\itemsep}{0.16em}%
  \setlength{\parsep}{0.1em}%
}{\end{enumerate}}
\newenvironment{examinstructions}{%
  \par\begingroup\itshape
}{%
  \par\endgroup\vspace{0.18em}
}
\makeatletter
\renewcommand\section{\@startsection{section}{1}{0pt}%
  {-1.25ex plus -.25ex minus -.1ex}%
  {0.55ex plus .1ex}%
  {\normalfont\large\bfseries}}
\renewcommand{\@maketitle}{%
  \newpage\null\vskip -1em%
  {\footnotesize\bfseries
    \href{https://www.ufl.edu/}{University of Florida}, \href{https://math.ufl.edu/}{Department of Mathematics}\par}%
  \vskip -0.5em%
  \tagstructbegin{tag=Title}%
    {\LARGE\bfseries\href{https://gma.math.ufl.edu/exams/analysis-phd/}{Analysis PhD Exam}, September 2000\par}%
  \tagstructend%
  \vskip 0.25em%
  \tagmcbegin{artifact}%
    \hrule height 1pt%
  \tagmcend%
  \vskip 1em%
}
\makeatother
\title{Analysis PhD Exam, September 2000}
\author{}
\date{}
\begin{document}
\maketitle
\thispagestyle{empty}
\begin{examinstructions}
Write all proofs in a neat and logical fashion.
\end{examinstructions}
\begin{examproblems}{0}
\item
State and prove the Radon–Nikodym Theorem.
\item
State and prove the Vitali Integral Convergence Theorem.
\item
Let \((X,\Sigma,\mu)\) be a finite measure space and suppose \(\Sigma=\sigma(\Sigma_0)\), where \(\Sigma_0\) is an algebra. Show that for every \(\varepsilon>0\), if \(E\in\Sigma\), then there exists an \(E_0\in\Sigma_0\) such that \(\mu(E\mathbin{\triangle}E_0)<\varepsilon\).
\item
Let~\(Y\) and~\(Z\) be closed subspaces in the B-space~\(X\). Suppose each \(x\in X\) has a unique representation in the form \(x=y+z\), with \(y\in Y\) and \(z\in Z\). Show that there exists a constant~\(K\) such that \(\lVert y\rVert\leq K\lVert x\rVert\) and \(\lVert z\rVert\leq K\lVert x\rVert\), for each \(x\in X\). [Hint: First apply the Closed Graph Theorem to \(T:X\to Y\), \(T(x)=y\).]
\item
Let \((X,\Sigma,\mu)\) be a measure space. Let \((f_n)\) be a sequence of integrable functions such that \(\sum_n\int |f_n|\,d\mu<\infty\). Does \(\sum_n f_n\) converge a.e.? Prove.
\item
Let \((\Omega,\mathcal F,P)\) be a probability space. Let~\(X\) and~\(Y\) be integrable random variables. Suppose~\(X\) and~\(Y\) are independent. Find \(E(X/Y)\).
\item
Let \((X,\Sigma,\mu)\) be a finite measure space; let \((f_n)\) be a sequence of integrable functions such that \(\int_E f_n\,d\mu\) converges for each \(E\in\Sigma\). Assume \((f_n)\) is \(L^1\)-bounded. Show there exists an integrable function~\(f\) such that \(\int_E f_n\,d\mu\to\int_E f\,d\mu\) for all \(E\in\Sigma\).
\end{examproblems}
\end{document}
